\subsection{例：单纯形锥与单纯形}

\begin{enumerate}
\def\labelenumi{\alph{enumi})}
\tightlist
\item
  设 \(a\) 是 \(\mathbf{E}\) 中的一点，\((e_1, \ldots, e_d)\) 是 \(\mathbf{T}\) 的一组基；于是 \(\mathbf{E}\) 中每一点都可唯一写成下列形式
\end{enumerate}

\[
x = a + \xi_ {1}. e _ {1} + \dots + \xi_ {d}. e _ {d}\tag{7}
\]

其中 \(\xi_1, \ldots, \xi_d\) 是实数。记 \(e_i'\) 为 E 上的仿射函数，它把任意写成形式 (7) 的 \(x \in \mathbf{E}\) 赋值 \(\xi_i\)。记 \(\mathbf{H}_i\) 为由满足 \(x\) 的那些 \(e_i'(x) = 0\) 组成的超平面，记 \(\mathfrak{H}\) 为超平面 \(\mathbf{H}_1, \ldots, \mathbf{H}_d\) 所成的集合。对于集合 I = \{1, 2, \ldots, d\} 的每个子集 J，置 \(\mathbf{H}_J = \bigcap_{i \in J} \mathbf{H}_i\)；对于由 0、1 或 -1 组成的每个数列 \((\varepsilon_1, \ldots, \varepsilon_d)\)，用 \(\mathbf{F}(\varepsilon_1, \ldots, \varepsilon_d)\) 表示所有 \(x \in \mathbf{E}\) 中满足 \(e_i'(x)\) 的符号为 \(\varepsilon_i\) 且对所有 \(i\) 属于 I 的那些点所成的集合（《一般拓扑》，第 4 章，§3，第 2 号）。立即可见，\(\mathfrak{H}\) 在 E 中定义的面片就是集合 \(\mathbf{F}(\varepsilon_1, \ldots, \varepsilon_d)\)，且这些集合两两不同；\(\mathbf{F}(\varepsilon_1, \ldots, \varepsilon_d)\) 的支撑等于 \(\mathbf{H}_J\)，其中 J 是由满足对于 \(i \in I\)，\(\varepsilon_i = 0\) 的那些 i 组成的集合；特别地，室就是形如 \(\mathbf{F}(\varepsilon_1, \ldots, \varepsilon_d)\) 且每个 \(\varepsilon_i\) 等于 1 或 -1 的集合。

集合 \(\mathrm{C} = \mathrm{F}(1, \ldots, 1)\)，即由那些 \(x\) 满足 \(e_i'(x) > 0\) 且对所有 \(i \in \mathrm{I}\) 的点组成，是一个室，称为在顶点 \(a\) 处由基 \((e_{1},\ldots,e_{d})\) 定义的开单纯形锥。其闭包由满足 \(e_{i}^{\prime}(x)\geqslant0\) 对所有 \(i\in I\) 且 \(d\geqslant1\) 的点 x 组成；否则，其闭包为空。对于 I 的每个子集 J，令 \(C_{J}\) 为 E 中满足 \(e_{i}^{\prime}(x)=0\) 对 \(i\in J\) 且 \(e_{i}^{\prime}(x)>0\) 对 \(i\in I-J\) 的点 x 所成的集合。于是 \(C_{J}\) 是支撑为 \(H_{J}\) 的面片，并且是仿射空间 \(H_{J}\) 中顶点为 a 的开单纯形锥；此外，\(\overline{C}=\bigcup_{J\subset I}C_{J}\)。特别地，C 的墙是超平面 \(H_{i}\)（\(i\in I\)），而包含于 \(H_{i}\) 的 C 的面等于 \(C_{\{i\}}\)。

集合 \(\mathrm{H}_i, \mathrm{H}_J, \mathrm{C}, \mathrm{C}_J\) 和 \(\mathrm{F}(\varepsilon_1, \ldots, \varepsilon_d)\) 在把基 \((e_1, \ldots, e_d)\) 换成基 \((\lambda_1 e_1, \ldots, \lambda_d e_d)\)（其中对所有 \(\lambda_i > 0\) 有 \(i\)）时均不改变。

\begin{enumerate}
\def\labelenumi{\alph{enumi})}
\setcounter{enumi}{1}
\tightlist
\item
  现在假设给定 E 的一个仿射自由点系，例如 \((a_{0}, a_{1}, \ldots, a_{d})\)。我们知道 E 的每一点都能唯一写成 \(x = \xi_{0}.a_{0} + \cdots + \xi_{d}.a_{d}\) 的形式，其中 \(\xi_{0}, \ldots, \xi_{d}\) 是满足 \(\xi_{0} + \cdots + \xi_{d} = 1\) 的实数（《代数》，第 2 章，§9，第 3 号）。定义仿射函数 \(f_{0}, \ldots, f_{d}\)，其中函数 \(f_{i}\) 把如上写出的点 x 赋值为数 \(\xi_{i}\)。记 \(H_{i}\) 为由 \(f_{i}(x) = 0\) 定义的 E 的超平面，记 H 为超平面 \(H_{0}, H_{1}, \ldots, H_{d}\) 所成的集合；最后置 \(I = \{0, 1, \ldots, d\}\)。以 \(a_{0}, \ldots, a_{d}\) 为顶点的开单纯形，是 E 中满足对所有 \(f_{i}(x) > 0\) 有 \(i \in I\) 的点 x 所成的集合 C；它是 H 在 E 中定义的一个室。C 的闭包是点 x 所成的集合 \(\overline{C}\)，其中对所有 \(f_{i}(x) \geqslant 0\) 有 \(i \in I\)；它是有限集 \(\{a_{0}, \ldots, a_{d}\}\) 的凸包，并且容易看出 \(\overline{C}\) 的极点是 \(a_{0}, \ldots, a_{d}\)。
\end{enumerate}

对于任意子集 \(J\)，若它是 \(I\) 的且不同于 \(I\)，置 \(H_J = \bigcap_{i \in J} H_i\)，并令 \(C_J\) 为点 \(x\) 所成的 \(E\) 的集合，满足 \(f_i(x) = 0\) 对 \(i \in J\) 且 \(f_i(x) > 0\) 对 \(i \in I - J\)。集合 \(C_J\) 是仿射空间 \(H_J\) 中以点 \(a_i\)（对 \(i \in I - J\)）为顶点的开单纯形；有 \(C_\emptyset = C, \overline{C} = \bigcup_{J \neq I} C_J\)，且 \(C_J \neq C_{J'}\) 当 \(J \neq J'\) 时；此外，\(C_J\) 是支撑为 \(H_J\) 的面片。特别地，\(C\) 的墙是 \(H_0, \ldots, H_d\)，而 \(C_{\{i\}}\) 是包含于 \(H_i\) 的面。

对于非空子集 \(K\)，若它是 \(I\) 的，令 \(B_K\) 为点 \(x\) 所成的 \(E\) 的集合，满足 \(f_i(x) > 0\) 对 \(i \in K\) 且 \(f_i(x) < 0\) 对 \(i \in I - K\)。集合 \(B_K\) 是由 \(\mathfrak{H}\) 在 \(E\) 中定义的室，并且 \(B_I = C\)。容易看出 \(\overline{C}\) 是紧的；另一方面，若 \(K\) 是 \(I\) 的一个不同于 I 的子集，基数为 \(p\)，则室 \(B_K\) 包含由下式定义的点序列 \(x_n\)，其中 \(n \geqslant 2\)：

\[
f _ {i} (x _ {n}) = \left\{ \begin{array}{l l} n & \text { for } i \in K \\ (1 - p n) / (d + 1 - p) & \text { for } i \in I - K \end{array} \right.
\]

这表明 \(\mathsf{B}_{\mathsf{K}}\) 不是相对紧的。

